\subsubsection{Piles au lithium}
\label{subsubsection:piles_au_lithium}
%Marucco exercice 3.2

\setcounter{equation}{0}
\setcounter{Q}{0}

Considérons les piles au lithium suivantes, dont l'électrolyte solide (ES) est l'iodure de lithium : \ce{LiI}

\begin{itemize}
\item Piles utilisées pour les stimulateurs cardiaques et les calculatrices de poche, fabriquées par centaines de milliers d'exemplaires :
\begin{equation*}
\ominus ~ \text{acier} | \ce{Li} | \ce{LiI (ES)} | \ce{I2 (P2VP)} | \text{acier} ~ \oplus
\end{equation*}

P2VP est le complexe I2-poly-2-vinyle-pyridine à 92-94 \si{\percent} d'iode.
\item Piles sans autodécharge pouvant être stockées pendant 2 ans :

\begin{equation*}
\ominus ~ \text{acier} | \ce{Li}  | \ce{LiI-Al2O3 (ES)} | \ce{PbI2}, \ce{PbS} | \text{acier} ~ \oplus
\end{equation*}

%\item Piles utilisées pour les calculatrices et dans les têtes de forages pétroliers :

%\begin{equation*}
%\ominus ~ \text{acier} | \ce{Li} | \ce{LiI-Li4P2S7 (ES)} | \ce{TiS2}  | \text{acier} ~ \oplus
%\end{equation*}
%On notera que \ce{TiS2} est un matériau d'insertion.

\end{itemize}

\paragraph*{\'{E}tude de la pile $\ce{Li} | \ce{LiI (ES)} | \ce{I2 (P2VP)}$}

\Q \label{Q:01-01a} Rappeler l'expression du potentiel d'électrode $E_i$ à l'interface $i$ en fonction des potentiels électrostatiques. En déduire la relation existant entre enthalpie libre et potentiels électrostatiques à chacune des interfaces.
\solution{Dans le cours, nous avons vu que le potentiel d'électrode représente la différence de potentiel électrique à l'interface entre l'électrode et la "solution" :
\begin{align*}
E_i = \Phi_i (\text{électrode}) - \Phi_i ( \text{solution} )
\end{align*}
\`{A} chacune des interfaces, à l'équilibre, on a : 
\begin{align*}
\Delta_{1/2} G_i &= - n F \left( \Phi_i (\text{électrode}) - \Phi_i ( \text{solution} ) \right) \\
&= -n F E_i
\end{align*}
où $n$ représente le nombre d'électrons échangé à l'interface par mole d'oxydant.
}

\Q \label{Q:01-02a} Quelles sont les réactions que l'on peut considérer aux interfaces des électrodes ? En déduire l'équation de réaction.
\solution{
\begin{center}
\begin{tabular}{ccc}
Interface & Réaction & Enthalpie libre de demi-réaction \\
$\ominus$ & $\ce{Li = Li^+ + e^-}$ & $\Delta_{1/2} G_\ominus$ \\
$\oplus$ & $\ce{I^- = \frac{1}{2} I_2 + e^-}$ & $\Delta_{1/2} G_\oplus$ \\
\end{tabular}
\end{center}
Au bilan, on trouve l'équation de réaction suivante : 
\begin{align*}
\ce{Li + \frac{1}{2} I2 = Li^+ + I^-} \implies \ce{Li (\ominus) + \frac{1}{2} I2  (\oplus) = LiI (\text{el})} \quad \Delta_r G
\end{align*}}

\Q \label{Q:01-03a} Exprimer les enthalpies de demi-pile à chacune des interfaces en fonction des potentiels chimiques des espèces, puis dans l'enthalpie libre de la réaction "fictive" en fonction de ces mêmes potentiels.
\solution{
Nous avons vu en cours que l'enthalpie de réaction s'exprime comme une combinaison linéaire des potentiels chimiques : 
\begin{align*}
\Delta_r G = \sum_i \nu_i \mu_i
\end{align*}
Par conséquent, on trouve : 
\begin{align*}
\Delta_{1/2} G_\ominus &= \mu_{el} ( \ce{Li^+} ) + \mu_\ominus ( \ce{e^-} ) - \mu_\ominus ( \ce{Li} ) \\
\Delta_{1/2} G_\oplus &= \frac{1}{2} \mu_\oplus ( \ce{I2} ) + \mu_\oplus ( \ce{e^-} ) - \mu_{el} ( \ce{I^-} ) \\
\Delta_r G &= \mu_{el} ( \ce{LiI} ) - \mu_\ominus ( \ce{Li} ) - \frac{1}{2} \mu_\oplus ( \ce{I2} ) 
\end{align*}
}

\Q \label{Q:01-04a} \'{E}tablir une relation entre les 3 relations de la Q. \ref{Q:01-03a} et la force électromotrice ou les potentiels des électrodes.
\solution{En utilisant la Q. \ref{Q:01-01a} \& \ref{Q:01-03a}
\begin{align*}
-F E_\ominus &= \mu_{el} ( \ce{Li^+} ) + \mu_\ominus ( \ce{e^-} ) - \mu_\ominus ( \ce{Li} ) \\
-F E_\oplus &= \frac{1}{2} \mu_\oplus ( \ce{I2} ) + \mu_\oplus ( \ce{e^-} ) - \mu_{el} ( \ce{I^-} ) \\
-F e &= \mu_{el} ( \ce{LiI} ) - \mu_\ominus ( \ce{Li} ) - \frac{1}{2} \mu_\oplus ( \ce{I2} ) 
\end{align*}
} 

\Q Montrer que la f.e.m s'exprime bien comme la différence des potentiels d'électrode.
\solution{
\begin{align*}
-F \left( E_{\oplus} - E_\ominus \right) &= 
\left( \frac{1}{2} \mu_\oplus ( \ce{I2} ) + \mu_\oplus ( \ce{e^-} ) - \mu_{el} ( \ce{I^-} ) \right) 
- \left( \mu_{el} ( \ce{Li^+} ) + \mu_\ominus ( \ce{e^-} ) - \mu_\ominus ( \ce{Li} ) \right) \\
&= \left( \frac{1}{2} \mu_\oplus ( \ce{I2} ) + \mu_\ominus ( \ce{Li} ) \right) - \left( \mu_{el} ( \ce{Li^+} + \mu_{el} ( \ce{I^-} ) \right) + \left( \mu_\oplus ( \ce{e^-} ) - \mu_\ominus ( \ce{e^-} ) \right) 
\end{align*}
Or, en cours, on a demontré la relation d'équilibre des phases qui impose $-Fe = \mu_\oplus ( \ce{e^-} ) - \mu_\ominus ( \ce{e^-} )$, par conséquent, on a $e = E_\oplus - E_\ominus$ si et seulement si : 
\begin{align*}
\left( \frac{1}{2} \mu_\oplus ( \ce{I2} ) + \mu_\ominus ( \ce{Li} ) \right) = \left( \mu_{el} ( \ce{Li^+} + \mu_{el} ( \ce{I^-} ) \right)
\end{align*}
Condition mathématique vérifiée puisqu'elle correspond à deux façons différentes de former \ce{LiI} (en partant des corps simples ou en partant des espèces ioniques).
}

\Q Exprimer l'enthalpie de formation standard de la phase nouvellement formée.
\solution{La phase formée est \ce{LiI} dont la réaction de formation est : 
\begin{align*}
\ce{Li + I2 = LiI} 
\end{align*}
L'enthalpie libre de formation standard s'exprime comme :
\begin{align*}
\Delta_f G^\circ ( \ce{LiI} ) = \mu_{el}^\circ ( \ce{LiI} ) - \frac{1}{2} \mu_\oplus^\circ ( \ce{I2} ) - \frac{1}{2} \mu_\ominus^\circ ( \ce{Li} )
\end{align*}
}

\Q En déduire la relation existant entre $\Delta_r G$, $\Delta_f G^\circ ( \ce{LiI} )$ et d'autres paramètres à déterminer. En vous plaçant dans l'idéalité, on assimile l'activité d'un soluté à sa fraction molaire, en déduire une bonne approximation entre $\Delta_r G$ et $\Delta_f G^\circ ( \ce{LiI} )$.
\solution{Sachant que Li et LiI sont deux solides, corps purs, leur activité vaut $1$, on a alors : 
\begin{align*}
\mu_\ominus ( \ce{Li} ) &= \mu_\ominus^\circ ( \ce{Li} ) \\
\mu_{el} ( \ce{LiI} ) &= \mu_{el}^\circ ( \ce{LiI} ) \\
\mu_\oplus ( \ce{I2} ) &= \mu_\oplus^\circ ( \ce{I2} ) + RT \ln ( a_\oplus (\ce{I2}) ) \\
\end{align*}
On trouve alors que : 
\begin{align*}
\Delta_r G^\circ &= \underbrace{\mu_{el}^\circ ( \ce{LiI} ) - \mu_\ominus^\circ ( \ce{Li} ) - \frac{1}{2} \mu_\oplus^\circ ( \ce{I2} )}_{\Delta_f G^\circ ( \ce{LiI} ) } + RT \ln ( a_\oplus (I2) ) \\
&= \Delta_f G^\circ ( \ce{LiI} ) + RT \ln ( a_\oplus ( \ce{I2} ) )
\end{align*}
On sait que le complexe formé entre le polymère et \ce{I2} contient $x ( \ce{I2} ) \approx 0.92 $. On s'attend à avoir $RT | \ln (x (\ce{I2})) | < 0.2 ~ \si{\kilo\joule\per\mole}$, ce qui est négligeable devant l'enthalpie de formation de \ce{LiI} (de l'ordre d'une énergie de liaison $> 100^{\text{aine}} ~ \si{\joule\per\mole}$).
Si cette approximation est vérifiée, on a : 
\begin{align}
\Delta_r G \approx \Delta_f G^\circ (LiI)
\end{align}}

En réalisant plusieurs titrages coulométriques à différentes températures, il a été montré que :
\begin{align*}
\Delta_r G^\circ (T) = -270864 - 85.7 T \quad [\si{\joule\per\mole}]
\end{align*}
avec $T$ en \si{\kelvin}.
Cette relation est vérifiée à température ambiante.

\Q En déduire une valeur de l'enthalpie de formation de \ce{LiI} à \SI{25}{\celsius}.
\solution{\begin{align*}
\Delta_f G^\circ ( 25 ~ \si{\celsius} ) \approx \Delta_r G ( 25 ~ \si{\celsius} ) \stackrel{AN}{\approx} -245 ~ \si{\kilo\joule\per\mole}
\end{align*}}

\paragraph*{\'{E}tude de la pile $\ce{Li}  | \ce{LiI-Al2O3 (ES)} | \ce{PbI2}, \ce{PbS}$}

\Q Quelles sont les réactions chimiques se produisant à cette pile et quelle est la grandeur thermodynamique que l'on peut déterminer ?
\solution{
Le pôle $\oplus$ est constitué de deux constituants pouvant réagir de manière rédox. On s'attend à avoir deux réactions rédox pour cette pile : 
\begin{center}
\begin{minipage}{0.48\textwidth}
\begin{align*}
\ce{2 Li &= 2 Li^+ + 2 e^-} \\ 
\ce{PbI2 + 2 e^- &= Pb + 2I^-} \\
\hline
\ce{PbI2 + 2 Li &= Pb + 2 Li^+ + 2 I^-}\\
\ce{PbI2 + 2 Li &= Pb + 2 LiI}
\end{align*}
\end{minipage}\hfill \begin{minipage}{0.48\textwidth}
\begin{align*}
\ce{2 Li &= 2 Li^+ + 2 e^-} \\ 
\ce{PbS + 2 e^- &= Pb + S^{2-}} \\
\hline
\ce{PbS + 2 Li &= Pb + 2Li^+ + S^{2-}} \\
\ce{PbS + 2 Li &= Pb + Li_2 S} 
\end{align*}
\end{minipage}
\end{center}
Traitons l'équation de gauche, on peut alors exprimer l'enthalpie de réaction : 
\begin{align}
\Delta_r G = \mu ( \ce{Pb} ) + 2\mu ( \ce{LiI} ) - \mu ( \ce{PbI2} ) - 2 \mu ( \ce{Li} ) 
\end{align}
Or, on peut exprimer l'enthalpie de formation de chacun des corps purs complexes : 
\begin{align*}
\ce{Li + \frac{1}{2} I2 &= LiI} \quad \Delta_f G^\circ ( \ce{LiI} ) = \mu^\circ ( \ce{LiI}) - \mu^\circ ( \ce{Li} ) - \frac{1}{2} \mu^\circ ( \ce{I2} )   \\
\ce{Pb + I2 &= PbI2} \quad \Delta_f G^\circ ( \ce{PbI2} ) = \mu^\circ ( \ce{PbI2}) - \mu^\circ ( \ce{I2} ) - \mu^\circ ( \ce{Pb} ) 
\end{align*}
Tous les corps formés sont des corps purs, on peut donc écrire : 
\begin{align*}
\Delta_r G &= \mu^\circ ( \ce{Pb} ) + \mu^\circ ( \ce{LiI} ) - \mu^\circ ( \ce{PbI2} ) - 2 \mu^\circ ( \ce{Li} )  \\
&= \mu^\circ ( \ce{Pb} ) + 2( \Delta_f G^\circ ( \ce{LiI} ) + \mu^\circ ( \ce{Li} ) + \frac{1}{2} \mu^\circ ( \ce{I2} ) ) -  ( \Delta_f G^\circ ( \ce{PbI2} ) + \mu^\circ ( \ce{I2} ) + \mu^\circ ( \ce{Pb} ) ) - 2 \mu^\circ ( \ce{Li} ) \\
&= 2 \Delta_f G^\circ ( \ce{LiI} ) - \Delta_f G^\circ ( \ce{PbI2} ) 
\end{align*}
La mesure de l'enthalpie de réaction permet de mesurer la stabilité relative ls phases formées \ce{LiI} et \ce{PbI2}}

%\paragraph*{\'{E}tude de la pile $\ce{Li} | \ce{LiI-Li4P2S7 (ES)} | \ce{TiS2}$}

%\Q Quelle est la réaction chimique bilan se produisant à cette pile et quelle est la grandeur thermodynamique que l'on peut déterminer ?